% Edit only the text between matching begin and end lines.
\documentclass[11pt,a4paper]{article}
% Shared preamble for TMUA question papers. Local edits here are overwritten.

% Reproducible PDFs: no timestamps, no trailer id, no build-path details.
\ifdefined\pdfsuppressptexinfo
  \pdfsuppressptexinfo=-1
  \pdfinfoomitdate=1
  \pdftrailerid{}
\fi

\usepackage{graphicx}
\usepackage{mathtools}
\usepackage{amssymb}
\usepackage{amsmath}
\usepackage{amsthm}
\usepackage{enumitem}
\usepackage{microtype}
\usepackage[table]{xcolor}
\usepackage{array}
\usepackage{booktabs}
\usepackage{tabularx}
\usepackage{longtable}
\usepackage{ragged2e}
\usepackage{fancyhdr}
\usepackage{etoolbox}
\usepackage[
  a4paper,
  left=0.8in,
  right=1in,
  top=1in,
  bottom=0.5in,
  headheight=14pt,
  headsep=0.4in
]{geometry}

\usepackage{pgfplots}
\pgfplotsset{compat=1.18}
\usepgfplotslibrary{fillbetween, polar}
\usetikzlibrary{calc, positioning}
\usetikzlibrary{angles, quotes}
\usetikzlibrary{arrows.meta}
\usetikzlibrary{decorations.pathreplacing, decorations.markings}
\usetikzlibrary{patterns}
\usetikzlibrary{intersections}

\usepackage{tocloft}
\usepackage[hidelinks]{hyperref}
\hypersetup{pdfauthor={danieljsmith.org}, pdfcreator={}, pdfsubject={}, pdfkeywords={}}
\urlstyle{same}
\tocloftpagestyle{djspaper}
\setlength{\cftbeforesubsecskip}{2pt}
\setlength{\cftsubsecindent}{1.5em}
\setlength{\cftsubsecnumwidth}{0pt}
\renewcommand{\cftsubsecfont}{\small\color{black}}
\renewcommand{\cftsubsecpagefont}{\small\color{black}}
\renewcommand{\cftsubsecleader}{\hfill}

% \djsPaperSetup{<PDF title>}{<header label>}, called once after this file.
\newcommand{\djsHeaderLabel}{}
\newcommand{\djsPaperSetup}[2]{%
  \hypersetup{pdftitle={#1}}%
  \renewcommand{\djsHeaderLabel}{#2}}

\setlength{\parindent}{0pt}
\fancypagestyle{djspaper}{%
  \fancyhead{}%
  \fancyfoot{}%
  \renewcommand{\headrulewidth}{0.9pt}%
  \fancyhead[L]{\djsHeaderLabel}%
  \fancyhead[R]{\href{https://danieljsmith.org}{danieljsmith.org}}%
}
\pagestyle{djspaper}

\newcolumntype{Y}{>{\RaggedRight\arraybackslash}X}

% Topic labels: a subtle line above each question and a suffix on its
% contents entry. The bookmark form is spelled out separately, because a PDF
% bookmark is plain text and would otherwise keep the colour name.
\newcommand{\djsTocTopic}[1]{%
  \ifblank{#1}{}{%
    \texorpdfstring
      {\hspace{0.9em}{\footnotesize\itshape\color{black!45}#1}}
      { \textendash\ #1}}}
\newcommand{\djsTopicLine}[1]{%
  \ifblank{#1}{}{%
    {\raggedleft\footnotesize\itshape\color{black!45}#1\par}%
    \vspace{-0.6\baselineskip}}}

\newcommand{\questionitem}[1][]{%
  \item\phantomsection
  \addcontentsline{toc}{subsection}{%
    Question \number\value{enumi}\protect\djsTocTopic{#1}}%
}
\renewcommand{\contentsname}{Questions}
\newcommand{\djsFrontMatter}{%
  \vspace*{20pt}%
  \tableofcontents
  \newpage
}

\djsPaperSetup{TMUA Set B Paper 2}{TMUA Set B - Paper 2}
\begin{document}
\thispagestyle{djspaper}
\djsFrontMatter
\thispagestyle{djspaper}
\vspace*{8pt}
\djsTopicLine{Polynomials}
\begin{enumerate}[label=\textbf{\arabic*.},leftmargin=*,itemsep=0pt,start=1]
\questionitem[Polynomials]
% begin q000071 stem
For a real constant $k$, consider the equation \[ x^4-4x^2+k=0 \] Which of the following conditions on $k$ is \textbf{necessary and sufficient} for the equation to have exactly two distinct real solutions?
% end q000071 stem

\par\vspace*{12pt}
\begin{tabularx}{\linewidth}{@{}>{\bfseries}l Y@{}}
A &
% begin q000071 choice A
$k>4$
% end q000071 choice A
\\[0.8em]
B &
% begin q000071 choice B
$k=4$
% end q000071 choice B
\\[0.8em]
C &
% begin q000071 choice C
$k\leq0$
% end q000071 choice C
\\[0.8em]
D &
% begin q000071 choice D
$0<k<4$
% end q000071 choice D
\\[0.8em]
E &
% begin q000071 choice E
$k=0$
% end q000071 choice E
\\[0.8em]
F &
% begin q000071 choice F
$k<0$ or $k=4$
% end q000071 choice F
\\[0.8em]
G &
% begin q000071 choice G
$k<0$
% end q000071 choice G
\\
\end{tabularx}
\end{enumerate}
\clearpage
\thispagestyle{djspaper}
\vspace*{8pt}
\djsTopicLine{Number}
\begin{enumerate}[label=\textbf{\arabic*.},leftmargin=*,itemsep=0pt,start=2]
\questionitem[Number]
% begin q000108 stem
Let $n$ be an integer. Consider the following conditions:
\[
\begin{array}{@{}l l@{}}
\textbf{I:} & \text{$n$ is odd} \\[2pt]
\textbf{II:} & \text{$n$ is a multiple of $4$} \\[2pt]
\textbf{III:} & \text{$n$ is a prime greater than $3$}
\end{array}
\]
Which of these conditions is/are \textbf{sufficient} to guarantee that $n^3-n$ is divisible by $12$?
% end q000108 stem

\par\vspace*{12pt}
\begin{tabularx}{\linewidth}{@{}>{\bfseries}l Y@{}}
A &
% begin q000108 choice A
None of them
% end q000108 choice A
\\[0.8em]
B &
% begin q000108 choice B
I only
% end q000108 choice B
\\[0.8em]
C &
% begin q000108 choice C
II only
% end q000108 choice C
\\[0.8em]
D &
% begin q000108 choice D
III only
% end q000108 choice D
\\[0.8em]
E &
% begin q000108 choice E
I and II only
% end q000108 choice E
\\[0.8em]
F &
% begin q000108 choice F
I and III only
% end q000108 choice F
\\[0.8em]
G &
% begin q000108 choice G
II and III only
% end q000108 choice G
\\[0.8em]
H &
% begin q000108 choice H
I, II, and III
% end q000108 choice H
\\
\end{tabularx}
\end{enumerate}
\clearpage
\thispagestyle{djspaper}
\vspace*{8pt}
\djsTopicLine{Polynomials}
\begin{enumerate}[label=\textbf{\arabic*.},leftmargin=*,itemsep=0pt,start=3]
\questionitem[Polynomials]
% begin q000612 stem
The curve $C$ has equation $y=x^3-3x^2$

A tangent to $C$ passes through the point $P(0,1)$. The point at which the tangent touches $C$ has $x$-coordinate $t$

Find the complete set of possible values of $t$
% end q000612 stem

\par\vspace*{12pt}
\begin{tabularx}{\linewidth}{@{}>{\bfseries}l Y@{}}
A &
% begin q000612 choice A
$t=-\frac12$ only
% end q000612 choice A
\\[0.8em]
B &
% begin q000612 choice B
$t=1$ only
% end q000612 choice B
\\[0.8em]
C &
% begin q000612 choice C
$t=-\frac12$ or $t=1$
% end q000612 choice C
\\[0.8em]
D &
% begin q000612 choice D
$t=-1$ or $t=1$
% end q000612 choice D
\\[0.8em]
E &
% begin q000612 choice E
$t=\frac12$ or $t=1$
% end q000612 choice E
\\[0.8em]
F &
% begin q000612 choice F
$t=-\frac12$ or $t=\frac32$
% end q000612 choice F
\\
\end{tabularx}
\end{enumerate}
\clearpage
\thispagestyle{djspaper}
\vspace*{8pt}
\djsTopicLine{Algebra}
\begin{enumerate}[label=\textbf{\arabic*.},leftmargin=*,itemsep=0pt,start=4]
\questionitem[Algebra]
% begin q000016 stem
A student is required to solve the equation \[ 3\sqrt{x+1}=x-1 \]

Their attempt is as follows:

\[ \begin{aligned} 3\sqrt{x+1}=x-1 &\xRightarrow{\textbf{(I)}} 9(x+1)=(x-1)^2\\[10pt] 9(x+1)=(x-1)^2 &\xRightarrow{\textbf{(II)}} 9x+9=x^2-2x+1\\[10pt] 9x+9=x^2-2x+1 &\xRightarrow{\textbf{(III)}} x^2-11x-8=0\\[10pt] x^2-11x-8=0 &\xRightarrow{\textbf{(IV)}} x=\dfrac{11\pm 3\sqrt{17}}{2} \\[10pt] \end{aligned} \] They then conclude that the solutions of the original equation are $x=\dfrac{11+3\sqrt{17}}{2}$ and $x=\dfrac{11-3\sqrt{17}}{2}$. Which one of the following is true?
% end q000016 stem

\par\vspace*{12pt}
\begin{tabularx}{\linewidth}{@{}>{\bfseries}l Y@{}}
A &
% begin q000016 choice A
The solution is correct.
% end q000016 choice A
\\[0.8em]
B &
% begin q000016 choice B
Only one of $x=\dfrac{11+3\sqrt{17}}{2}$ and $x=\dfrac{11-3\sqrt{17}}{2}$ is correct and the error arises as a result of step (I).
% end q000016 choice B
\\[0.8em]
C &
% begin q000016 choice C
Only one of $x=\dfrac{11+3\sqrt{17}}{2}$ and $x=\dfrac{11-3\sqrt{17}}{2}$ is correct and the error arises as a result of step (II).
% end q000016 choice C
\\[0.8em]
D &
% begin q000016 choice D
Only one of $x=\dfrac{11+3\sqrt{17}}{2}$ and $x=\dfrac{11-3\sqrt{17}}{2}$ is correct and the error arises as a result of step (III).
% end q000016 choice D
\\[0.8em]
E &
% begin q000016 choice E
There is another value of $x$ that satisfies the original equation and the error arises as a result of step (I).
% end q000016 choice E
\\[0.8em]
F &
% begin q000016 choice F
There is another value of $x$ that satisfies the original equation and the error arises as a result of step (II).
% end q000016 choice F
\\[0.8em]
G &
% begin q000016 choice G
There is another value of $x$ that satisfies the original equation and the error arises as a result of step (III).
% end q000016 choice G
\\
\end{tabularx}
\end{enumerate}
\clearpage
\thispagestyle{djspaper}
\vspace*{8pt}
\djsTopicLine{Trigonometry}
\begin{enumerate}[label=\textbf{\arabic*.},leftmargin=*,itemsep=0pt,start=5]
\questionitem[Trigonometry]
% begin q000063 stem
Let $p$ be a real constant. Consider the equation \[\sin x \cos x = p\] on $0 \le x < 2\pi$

Which of the following statements is/are true? \[ \begin{array}{@{}l l@{}}
\textbf{I:} & \text{If } |p|<\tfrac12 \text{ then the equation has exactly 4 distinct solutions} \\[5pt]
\textbf{II:} & \text{If the equation has exactly 2 distinct solutions, then } |p|=\dfrac12
\end{array} \]
% end q000063 stem

\par\vspace*{12pt}
\begin{tabularx}{\linewidth}{@{}>{\bfseries}l Y@{}}
A &
% begin q000063 choice A
None of them
% end q000063 choice A
\\[0.8em]
B &
% begin q000063 choice B
I only
% end q000063 choice B
\\[0.8em]
C &
% begin q000063 choice C
II only
% end q000063 choice C
\\[0.8em]
D &
% begin q000063 choice D
I and II
% end q000063 choice D
\\
\end{tabularx}
\end{enumerate}
\clearpage
\thispagestyle{djspaper}
\vspace*{8pt}
\djsTopicLine{Logic}
\begin{enumerate}[label=\textbf{\arabic*.},leftmargin=*,itemsep=0pt,start=6]
\questionitem[Logic]
% begin q000148 stem
We say a sequence of real numbers $x_n$ is \textit{a Cauchy sequence} if
\[
\begin{array}{c}
\text{for all $\varepsilon>0$ there exists a positive integer $N \in\mathbb{N}$ such that} \\[2pt]
\text{for all $n,m \geq N$ we have $|x_n - x_m | < \varepsilon$}
\end{array}
\]
Which one of the following statements is true if and only if $x_n$ is \textbf{not} a Cauchy sequence?
% end q000148 stem

\par\vspace*{12pt}
\begin{tabularx}{\linewidth}{@{}>{\bfseries}l Y@{}}
A &
% begin q000148 choice A
For every $\varepsilon>0$ and every $N\in\mathbb{N}$ there exists $m,n\ge N$ with $|x_m-x_n|<\varepsilon$
% end q000148 choice A
\\[0.8em]
B &
% begin q000148 choice B
There exists $\varepsilon>0$ such that for every $N\in\mathbb{N}$ and for all $m,n\le N$ we have $|x_m-x_n|\ge \varepsilon$
% end q000148 choice B
\\[0.8em]
C &
% begin q000148 choice C
For every $\varepsilon>0$ there exists $N\in\mathbb{N}$ and there exists $m,n\ge N$ with $|x_m-x_n|\ge \varepsilon$
% end q000148 choice C
\\[0.8em]
D &
% begin q000148 choice D
For every $N\in\mathbb{N}$ there exists $\varepsilon>0$ such that for all $m,n\ge N$ we have $|x_m-x_n|\ge \varepsilon$
% end q000148 choice D
\\[0.8em]
E &
% begin q000148 choice E
There exists $\varepsilon>0$ such that for every $N\in\mathbb{N}$ there exist $m,n\ge N$ with $|x_m-x_n|\ge \varepsilon$
% end q000148 choice E
\\[0.8em]
F &
% begin q000148 choice F
There exists $\varepsilon>0$ such that for every $N\in\mathbb{N}$ there exist $m,n\ge N$ with $|x_m-x_n|<\varepsilon$
% end q000148 choice F
\\[0.8em]
G &
% begin q000148 choice G
There exists $N\in\mathbb{N}$ such that for all $\varepsilon>0$ and all $m,n\ge N$ we have $|x_m-x_n|\ge \varepsilon$
% end q000148 choice G
\\[0.8em]
H &
% begin q000148 choice H
There exists $\varepsilon>0$ and there exists $N\in\mathbb{N}$ such that for all $m,n\ge N$ we have $|x_m-x_n|\ge \varepsilon$
% end q000148 choice H
\\
\end{tabularx}
\end{enumerate}
\clearpage
\thispagestyle{djspaper}
\vspace*{8pt}
\djsTopicLine{Integration}
\begin{enumerate}[label=\textbf{\arabic*.},leftmargin=*,itemsep=0pt,start=7]
\questionitem[Integration]
% begin q000134 stem
For a real constant $k$, define \[ A(k)=\int_{-1}^{2} |x-k|\,\mathrm{d}x \] For which value of $k$ is $A(k)$ minimised?
% end q000134 stem

\par\vspace*{12pt}
\begin{tabularx}{\linewidth}{@{}>{\bfseries}l Y@{}}
A &
% begin q000134 choice A
$-1$
% end q000134 choice A
\\[0.8em]
B &
% begin q000134 choice B
$\dfrac{1}{2}$
% end q000134 choice B
\\[0.8em]
C &
% begin q000134 choice C
$1$
% end q000134 choice C
\\[0.8em]
D &
% begin q000134 choice D
$\dfrac{3}{2}$
% end q000134 choice D
\\[0.8em]
E &
% begin q000134 choice E
$2$
% end q000134 choice E
\\[0.8em]
F &
% begin q000134 choice F
$0$
% end q000134 choice F
\\
\end{tabularx}
\end{enumerate}
\clearpage
\thispagestyle{djspaper}
\vspace*{8pt}
\djsTopicLine{Proof \& Error Analysis}
\begin{enumerate}[label=\textbf{\arabic*.},leftmargin=*,itemsep=0pt,start=8]
\questionitem[Proof \& Error Analysis]
% begin q000015 stem
Consider the following problem for an integer $n$: \[ \left( \dfrac{1}{3} \right)^{n+2} \ge \left( \dfrac{1}{27} \right)^{n-1} \] A student produces the following argument:

\[ \begin{array}{@{}l l@{}}
\textbf{I:} & \text{$\left( \dfrac{1}{3} \right)^{n+2} \ge \left( \dfrac{1}{27} \right)^{n-1}$} \\[20pt]

\textbf{II:} & \text{$\log_{\frac{1}{3}} \left( \left( \dfrac{1}{3} \right)^{n+2} \right) \ge \log_{\frac{1}{3}} \left( \left( \dfrac{1}{27} \right)^{n-1} \right)$} \\[20pt]

\textbf{III:} & \text{$n+2 \ge (n-1)\log_{\frac{1}{3}} \left( \dfrac{1}{27} \right)$} \\[20pt]

\textbf{IV:} & \text{$n+2 \ge 3(n-1)$} \\[20pt]

\textbf{V:} & \text{$n+2 \ge 3n-3$} \\[20pt]

\textbf{VI:} & \text{$n \le 2$} \\[20pt]
\end{array} \] Which step (if any) in the argument is invalid?
% end q000015 stem

\par\vspace*{12pt}
\begin{tabularx}{\linewidth}{@{}>{\bfseries}l Y@{}}
A &
% begin q000015 choice A
There are no invalid steps; the argument is correct.
% end q000015 choice A
\\[0.8em]
B &
% begin q000015 choice B
Only step \textbf{I} is invalid; the rest are correct.
% end q000015 choice B
\\[0.8em]
C &
% begin q000015 choice C
Only step \textbf{II} is invalid; the rest are correct.
% end q000015 choice C
\\[0.8em]
D &
% begin q000015 choice D
Only step \textbf{III} is invalid; the rest are correct.
% end q000015 choice D
\\[0.8em]
E &
% begin q000015 choice E
Only step \textbf{IV} is invalid; the rest are correct.
% end q000015 choice E
\\[0.8em]
F &
% begin q000015 choice F
Only step \textbf{V} is invalid; the rest are correct.
% end q000015 choice F
\\[0.8em]
G &
% begin q000015 choice G
Only step \textbf{VI} is invalid; the rest are correct.
% end q000015 choice G
\\
\end{tabularx}
\end{enumerate}
\clearpage
\thispagestyle{djspaper}
\vspace*{8pt}
\djsTopicLine{Trigonometry}
\begin{enumerate}[label=\textbf{\arabic*.},leftmargin=*,itemsep=0pt,start=9]
\questionitem[Trigonometry]
% begin q000008 stem
The equation \[ 2\cos^2 x - k\sin x = 1 \] has $n$ distinct solutions in the range $0 \le x \le 2\pi$

Consider the following statements: \[ \begin{array}{@{}l l@{}}
\textbf{I:} & \text{$k = 0$ is \textbf{necessary} for $n = 2$} \\[2pt]
\textbf{II:} & \text{$k > 2$ \textbf{only if} $n = 0$}
\end{array} \]
Which of the following statements is/are true?
% end q000008 stem

\par\vspace*{12pt}
\begin{tabularx}{\linewidth}{@{}>{\bfseries}l Y@{}}
A &
% begin q000008 choice A
Neither of them
% end q000008 choice A
\\[0.8em]
B &
% begin q000008 choice B
I only
% end q000008 choice B
\\[0.8em]
C &
% begin q000008 choice C
II only
% end q000008 choice C
\\[0.8em]
D &
% begin q000008 choice D
I and II
% end q000008 choice D
\\
\end{tabularx}
\end{enumerate}
\clearpage
\thispagestyle{djspaper}
\vspace*{8pt}
\djsTopicLine{Inequalities}
\begin{enumerate}[label=\textbf{\arabic*.},leftmargin=*,itemsep=0pt,start=10]
\questionitem[Inequalities]
% begin q000020 stem
A real number $k$ has the property that \[ x^2+kx+1\ge 0 \] for all real $x$ satisfying $x^2-4x+3\le 0$

Which of the following describes the set of all such values of $k$?
% end q000020 stem

\par\vspace*{12pt}
\begin{tabularx}{\linewidth}{@{}>{\bfseries}l Y@{}}
A &
% begin q000020 choice A
$k\ge -2$
% end q000020 choice A
\\[0.8em]
B &
% begin q000020 choice B
$k\ge -\dfrac{10}{3}$
% end q000020 choice B
\\[0.8em]
C &
% begin q000020 choice C
$-2\le k\le 2$
% end q000020 choice C
\\[0.8em]
D &
% begin q000020 choice D
$-4\le k\le 2$
% end q000020 choice D
\\[0.8em]
E &
% begin q000020 choice E
$-4\le k\le 4$
% end q000020 choice E
\\[0.8em]
F &
% begin q000020 choice F
$-6\le k\le -2$
% end q000020 choice F
\\[0.8em]
G &
% begin q000020 choice G
$k\ge 0$
% end q000020 choice G
\\
\end{tabularx}
\end{enumerate}
\clearpage
\thispagestyle{djspaper}
\vspace*{8pt}
\djsTopicLine{Functions \& Graphs}
\begin{enumerate}[label=\textbf{\arabic*.},leftmargin=*,itemsep=0pt,start=11]
\questionitem[Functions \& Graphs]
% begin q000133 stem
Consider the curve defined by \[ (|x|-4)^2+(y-1)^2=1 \] How many distinct straight lines passing through the origin are tangents to this curve?
% end q000133 stem

\par\vspace*{12pt}
\begin{tabularx}{\linewidth}{@{}>{\bfseries}l Y@{}}
A &
% begin q000133 choice A
$0$
% end q000133 choice A
\\[0.8em]
B &
% begin q000133 choice B
$1$
% end q000133 choice B
\\[0.8em]
C &
% begin q000133 choice C
$2$
% end q000133 choice C
\\[0.8em]
D &
% begin q000133 choice D
$3$
% end q000133 choice D
\\[0.8em]
E &
% begin q000133 choice E
$4$
% end q000133 choice E
\\[0.8em]
F &
% begin q000133 choice F
$5$
% end q000133 choice F
\\
\end{tabularx}
\end{enumerate}
\clearpage
\thispagestyle{djspaper}
\vspace*{8pt}
\djsTopicLine{Proof \& Error Analysis}
\begin{enumerate}[label=\textbf{\arabic*.},leftmargin=*,itemsep=0pt,start=12]
\questionitem[Proof \& Error Analysis]
% begin q000647 stem
Let $a$ and $b$ be positive real numbers, neither of which is equal to $1$. 

A student attempts to prove the following claim:
\begin{center} \begin{minipage}{0.82\linewidth} \centering \textbf{Claim:} If  $\log_a b=\log_b a$ then $a=b$ \end{minipage} \end{center}
The student writes the following argument: \[
\begin{array}{@{}ll}
\textbf{Line 1:} & \text{Let } t=\log_a b,\text{ so that }a^t=b \\[4pt]

\textbf{Line 2:} & \text{Since }b\neq1,\text{ we have }t\neq0 \\[4pt]

\textbf{Line 3:} &
\text{Raising both sides of }a^t=b\text{ to the power }\frac1t
\text{ gives }a=b^{1/t} \\[4pt]

\textbf{Line 4:} &
\text{Hence }\log_b a=\frac1t,\text{ so the given condition implies }
t=\frac1t,\text{ and therefore }t^2=1 \\[4pt]

\textbf{Line 5:} &
\text{Since }a^t=b\text{ and both }a\text{ and }b\text{ are positive, we must have }t>0 \\[4pt]

\textbf{Line 6:} & \text{Hence }t=1 \\[4pt]

\textbf{Line 7:} & \text{Therefore }b=a^1=a,\text{ which proves the claim.}
\end{array}
\]

Which line contains the first error?
% end q000647 stem

\par\vspace*{12pt}
\begin{tabularx}{\linewidth}{@{}>{\bfseries}l Y@{}}
A &
% begin q000647 choice A
Line 1
% end q000647 choice A
\\[0.8em]
B &
% begin q000647 choice B
Line 2
% end q000647 choice B
\\[0.8em]
C &
% begin q000647 choice C
Line 3
% end q000647 choice C
\\[0.8em]
D &
% begin q000647 choice D
Line 4
% end q000647 choice D
\\[0.8em]
E &
% begin q000647 choice E
Line 5
% end q000647 choice E
\\[0.8em]
F &
% begin q000647 choice F
Line 6
% end q000647 choice F
\\[0.8em]
G &
% begin q000647 choice G
Line 7
% end q000647 choice G
\\[0.8em]
H &
% begin q000647 choice H
There is no error; the argument is correct.
% end q000647 choice H
\\
\end{tabularx}
\end{enumerate}
\clearpage
\thispagestyle{djspaper}
\vspace*{8pt}
\djsTopicLine{Geometry}
\begin{enumerate}[label=\textbf{\arabic*.},leftmargin=*,itemsep=0pt,start=13]
\questionitem[Geometry]
% begin q000146 stem
Let $C$ be the circle $x^2+y^2=9$

Let $P$ be the point $(a,0)$ where $a>3$

The two tangents from $P$ to $C$ touch the circle at $T_1$ and $T_2$

What is the area of triangle $PT_1T_2$ in terms of $a$?
% end q000146 stem

\par\vspace*{12pt}
\begin{tabularx}{\linewidth}{@{}>{\bfseries}l Y@{}}
A &
% begin q000146 choice A
$3\sqrt{a^2-9}$
% end q000146 choice A
\\[0.8em]
B &
% begin q000146 choice B
$\dfrac{3(a^2-9)^{3/2}}{a^2}$
% end q000146 choice B
\\[0.8em]
C &
% begin q000146 choice C
$\dfrac{9}{a}\sqrt{a^2-9}$
% end q000146 choice C
\\[0.8em]
D &
% begin q000146 choice D
$\dfrac{a}{3}\sqrt{a^2-9}$
% end q000146 choice D
\\[0.8em]
E &
% begin q000146 choice E
$\dfrac{a^2-9}{3}$
% end q000146 choice E
\\[0.8em]
F &
% begin q000146 choice F
$\dfrac{a^2-9}{a}$
% end q000146 choice F
\\[0.8em]
G &
% begin q000146 choice G
$\dfrac{9}{\sqrt{a^2-9}}$
% end q000146 choice G
\\
\end{tabularx}
\end{enumerate}
\clearpage
\thispagestyle{djspaper}
\vspace*{8pt}
\djsTopicLine{Functions \& Graphs}
\begin{enumerate}[label=\textbf{\arabic*.},leftmargin=*,itemsep=0pt,start=14]
\questionitem[Functions \& Graphs]
% begin q000038 stem
Let $f(x)=|x-1|$ and $g(x)=|x+1|$

Define \[ h(x)=f\bigl(g(x)\bigr) \] How many integers $x$ satisfy $h(x)<1$?
% end q000038 stem

\par\vspace*{12pt}
\begin{tabularx}{\linewidth}{@{}>{\bfseries}l Y@{}}
A &
% begin q000038 choice A
$0$
% end q000038 choice A
\\[0.8em]
B &
% begin q000038 choice B
$1$
% end q000038 choice B
\\[0.8em]
C &
% begin q000038 choice C
$2$
% end q000038 choice C
\\[0.8em]
D &
% begin q000038 choice D
$3$
% end q000038 choice D
\\[0.8em]
E &
% begin q000038 choice E
$4$
% end q000038 choice E
\\[0.8em]
F &
% begin q000038 choice F
$5$
% end q000038 choice F
\\
\end{tabularx}
\end{enumerate}
\clearpage
\thispagestyle{djspaper}
\vspace*{8pt}
\djsTopicLine{Geometry}
\begin{enumerate}[label=\textbf{\arabic*.},leftmargin=*,itemsep=0pt,start=15]
\questionitem[Geometry]
% begin q000579 stem
A cuboid has edge lengths $a$, $b$ and $c$. The diagonals of its three different rectangular faces have lengths $5\text{ cm}$, $\sqrt{153}\text{ cm}$ and $4\sqrt{10}\text{ cm}$.

What is the volume of the cuboid?
% end q000579 stem

\par\vspace*{12pt}
\begin{tabularx}{\linewidth}{@{}>{\bfseries}l Y@{}}
A &
% begin q000579 choice A
$60\text{ cm}^3$
% end q000579 choice A
\\[0.8em]
B &
% begin q000579 choice B
$144\text{ cm}^3$
% end q000579 choice B
\\[0.8em]
C &
% begin q000579 choice C
$156\text{ cm}^3$
% end q000579 choice C
\\[0.8em]
D &
% begin q000579 choice D
$180\text{ cm}^3$
% end q000579 choice D
\\[0.8em]
E &
% begin q000579 choice E
$240\text{ cm}^3$
% end q000579 choice E
\\[0.8em]
F &
% begin q000579 choice F
$624\text{ cm}^3$
% end q000579 choice F
\\
\end{tabularx}
\end{enumerate}
\clearpage
\thispagestyle{djspaper}
\vspace*{8pt}
\djsTopicLine{Proof \& Error Analysis}
\begin{enumerate}[label=\textbf{\arabic*.},leftmargin=*,itemsep=0pt,start=16]
\questionitem[Proof \& Error Analysis]
% begin q000677 stem
A student attempts to prove the following claim.

\begin{center}
\begin{minipage}{0.82\linewidth}
\centering
\textbf{Claim:} For every real number $k$ with $k>\frac{2}{3}$ the equation \\ \qquad\qquad$x^3-3k^2x+k=0$ has three distinct real solutions.
\end{minipage}
\end{center}

The student produces the following argument:\\
\[
\begin{array}{@{}rl@{}}
\textbf{L1} & \text{Let } f(x)=x^3-3k^2x+k \text{, so } f'(x)=3x^2-3k^2 \text{, and } f'(x)=0 \text{ exactly when } \\ & x=-k \text{ or } x=k \text{,} 
  \text{and these two values are different because } k\neq0 \\[5pt]
\textbf{L2} & \text{We have } f'(x)>0 \text{ for } x<-k \text{ and for } x>k \text{, and } f'(x)<0 \text{ for } -k<x<k \\ & \text{so } f \text{ has a local maximum at } x=-k \text{ and a local minimum at } x=k \\[5pt]
\textbf{L3} & \text{The local maximum value is } f(-k)=2k^3+k \text{, which is positive because } k>0 \\[5pt]
\textbf{L4} & \text{The local minimum value is }
f(k)=k-2k^3<0\\[5pt]
\textbf{L5} & \text{A cubic with positive } x^3 \text{ coefficient whose local maximum value is positive and} \\
 & \text{whose local minimum value is negative meets the } x \text{-axis at three different points, so the} \\
 & \text{equation has three distinct real solutions}\\[5pt]
\end{array}
\]

Which one of the following statements about this argument is correct?
% end q000677 stem

\par\vspace*{12pt}
\begin{tabularx}{\linewidth}{@{}>{\bfseries}l Y@{}}
A &
% begin q000677 choice A
The first error is on line L1
% end q000677 choice A
\\[0.8em]
B &
% begin q000677 choice B
The first error is on line L2
% end q000677 choice B
\\[0.8em]
C &
% begin q000677 choice C
The first error is on line L3
% end q000677 choice C
\\[0.8em]
D &
% begin q000677 choice D
The first error is on line L4
% end q000677 choice D
\\[0.8em]
E &
% begin q000677 choice E
The first error is on line L5
% end q000677 choice E
\\[0.8em]
F &
% begin q000677 choice F
There is no error, and the claim is true.
% end q000677 choice F
\\[0.8em]
G &
% begin q000677 choice G
There is no error in any line, but the claim is false.
% end q000677 choice G
\\
\end{tabularx}
\end{enumerate}
\clearpage
\thispagestyle{djspaper}
\vspace*{8pt}
\djsTopicLine{Polynomials}
\begin{enumerate}[label=\textbf{\arabic*.},leftmargin=*,itemsep=0pt,start=17]
\questionitem[Polynomials]
% begin q000074 stem
Consider the statement:

\[ \begin{array}{c}
\text{For all real numbers $x$, \textbf{if} $x^2+px+1\le 0$ \textbf{then} $x\le 2$}
\end{array} \tag{$\ast$} \]

What is the complete set of real values of $p$ for which ($\ast$) is true?
% end q000074 stem

\par\vspace*{12pt}
\begin{tabularx}{\linewidth}{@{}>{\bfseries}l Y@{}}
A &
% begin q000074 choice A
no real values of $p$
% end q000074 choice A
\\[0.8em]
B &
% begin q000074 choice B
all real values of $p$
% end q000074 choice B
\\[0.8em]
C &
% begin q000074 choice C
$p\ge 0$
% end q000074 choice C
\\[0.8em]
D &
% begin q000074 choice D
$p<0$
% end q000074 choice D
\\[0.8em]
E &
% begin q000074 choice E
$p\ge -\dfrac{5}{2}$
% end q000074 choice E
\\[0.8em]
F &
% begin q000074 choice F
$p<-4$ \textbf{or} $p>0$
% end q000074 choice F
\\
\end{tabularx}
\end{enumerate}
\clearpage
\thispagestyle{djspaper}
\vspace*{8pt}
\djsTopicLine{Sequences \& Series}
\begin{enumerate}[label=\textbf{\arabic*.},leftmargin=*,itemsep=0pt,start=18]
\questionitem[Sequences \& Series]
% begin q000053 stem
A sequence $x_n$ is defined by $x_1=1$ and \[x_{n+1}=\frac{1}{3}x_n+2\] for $n\ge 1$.

Let $S$ be the sum to infinity of the series \[ \sum_{n=1}^{\infty} \left(x_{n+1}-x_n\right) \] What is the value of $S$?
% end q000053 stem

\par\vspace*{12pt}
\begin{tabularx}{\linewidth}{@{}>{\bfseries}l Y@{}}
A &
% begin q000053 choice A
$-\,\frac{3}{2}$
% end q000053 choice A
\\[0.8em]
B &
% begin q000053 choice B
$-1$
% end q000053 choice B
\\[0.8em]
C &
% begin q000053 choice C
$2$
% end q000053 choice C
\\[0.8em]
D &
% begin q000053 choice D
$\frac{9}{2}$
% end q000053 choice D
\\[0.8em]
E &
% begin q000053 choice E
The series does not converge.
% end q000053 choice E
\\
\end{tabularx}
\end{enumerate}
\clearpage
\thispagestyle{djspaper}
\vspace*{8pt}
\djsTopicLine{Number}
\begin{enumerate}[label=\textbf{\arabic*.},leftmargin=*,itemsep=0pt,start=19]
\questionitem[Number]
% begin q000578 stem
The highest common factor of two positive integers is $18$, and their lowest common multiple is $3780$.

What is the smallest possible value of the sum of the two integers?
% end q000578 stem

\par\vspace*{12pt}
\begin{tabularx}{\linewidth}{@{}>{\bfseries}l Y@{}}
A &
% begin q000578 choice A
$504$
% end q000578 choice A
\\[0.8em]
B &
% begin q000578 choice B
$522$
% end q000578 choice B
\\[0.8em]
C &
% begin q000578 choice C
$540$
% end q000578 choice C
\\[0.8em]
D &
% begin q000578 choice D
$558$
% end q000578 choice D
\\[0.8em]
E &
% begin q000578 choice E
$666$
% end q000578 choice E
\\[0.8em]
F &
% begin q000578 choice F
$738$
% end q000578 choice F
\\
\end{tabularx}
\end{enumerate}
\clearpage
\thispagestyle{djspaper}
\vspace*{8pt}
\djsTopicLine{Logic}
\begin{enumerate}[label=\textbf{\arabic*.},leftmargin=*,itemsep=0pt,start=20]
\questionitem[Logic]
% begin q000091 stem
Consider the statement about a real number $k$:

\[ \begin{array}{c}
\text{For all real numbers $x$ there exists a real number $y$ such that} \\ \text{ $(x-y)^2+k(x+y)>0$}
\end{array} \tag{$\ast$} \]

For which values of $k$ is ($\ast$) true?
% end q000091 stem

\par\vspace*{12pt}
\begin{tabularx}{\linewidth}{@{}>{\bfseries}l Y@{}}
A &
% begin q000091 choice A
for no real values of $k$
% end q000091 choice A
\\[0.8em]
B &
% begin q000091 choice B
for all real values of $k$
% end q000091 choice B
\\[0.8em]
C &
% begin q000091 choice C
for all real $k$ with $k>0$ 
% end q000091 choice C
\\[0.8em]
D &
% begin q000091 choice D
for all real $k$ with $k<0$ 
% end q000091 choice D
\\[0.8em]
E &
% begin q000091 choice E
for all real $k$ with $k\ne 0$ 
% end q000091 choice E
\\[0.8em]
F &
% begin q000091 choice F
for exactly one real value of $k$
% end q000091 choice F
\\[0.8em]
G &
% begin q000091 choice G
for all real $k$ with $k\geq0$ 
% end q000091 choice G
\\[0.8em]
H &
% begin q000091 choice H
for all real $k$ with $k\leq 0$ 
% end q000091 choice H
\\
\end{tabularx}
\end{enumerate}
\clearpage
\thispagestyle{djspaper}
\vspace*{8pt}
\phantomsection
\addcontentsline{toc}{section}{Answer Key}
{\large\bfseries Answer Key\par}
\vspace{8pt}
\begingroup
\renewcommand{\arraystretch}{1.15}
\setlength{\LTleft}{0pt}
\setlength{\LTright}{\fill}
\begin{longtable}{|c|c|}
\hline
\rowcolor{black!10}\textbf{Question} & \textbf{Answer} \\ \hline
\endfirsthead
\hline
\rowcolor{black!10}\textbf{Question} & \textbf{Answer} \\ \hline
\endhead
1 & F \\ \hline
2 & H \\ \hline
3 & C \\ \hline
4 & B \\ \hline
5 & D \\ \hline
6 & E \\ \hline
7 & B \\ \hline
8 & C \\ \hline
9 & A \\ \hline
10 & A \\ \hline
11 & D \\ \hline
12 & E \\ \hline
13 & B \\ \hline
14 & C \\ \hline
15 & B \\ \hline
16 & D \\ \hline
17 & E \\ \hline
18 & C \\ \hline
19 & B \\ \hline
20 & B \\ \hline
\end{longtable}
\endgroup
\end{document}
